\chapter*{Danh mục thuật ngữ và từ viết tắt}
\addcontentsline{toc}{chapter}{Danh mục thuật ngữ và từ viết tắt}

{\cobang
\begin{longtable}{P{3.6cm}P{11.6cm}}
\toprule
\textbf{Thuật ngữ} & \textbf{Ý nghĩa trong báo cáo}\\
\midrule
\endhead
Mô hình ngôn ngữ lớn (LLM) & Mô hình học máy được huấn luyện trên lượng văn bản rất lớn, có thể đọc câu hỏi và viết câu trả lời bằng ngôn ngữ tự nhiên.\\
Tác tử (agent) & Một thành phần phần mềm dùng mô hình ngôn ngữ để tự quyết định bước tiếp theo, ví dụ gọi công cụ nào, rồi dùng kết quả để trả lời.\\
SA & Một tác tử (Single-Agent): một tác tử duy nhất dùng toàn bộ công cụ.\\
MA & Nhiều tác tử (Multi-Agent): một bộ điều phối chia việc cho nhiều chuyên gia.\\
MA cố định & Cấu hình nhiều tác tử, kế hoạch do quy tắc định sẵn của chương trình lập ra và không chỉnh giữa chừng.\\
MA thích ứng & Cấu hình nhiều tác tử, kế hoạch do mô hình đề xuất và được phép bổ sung một lần sau khi có kết quả ban đầu.\\
MA không tra cứu & Như MA thích ứng nhưng tắt việc tra cứu kho tri thức.\\
Chuyên gia & Một thành phần chỉ lo một loại việc (sản phẩm, đánh giá, độ tin cậy, thị trường).\\
Bộ điều phối & Thành phần nhận câu hỏi, lập kế hoạch và giao việc cho các chuyên gia.\\
Tra cứu tài liệu (RAG) & Sinh câu trả lời có tra cứu (Retrieval-Augmented Generation): tìm các đoạn tài liệu liên quan trước, rồi mới để mô hình viết câu trả lời dựa trên các đoạn đó.\\
Kho tri thức & Tập tài liệu ngắn do nhóm tự viết lại từ nguồn chính thống (nhà xuất bản, thư viện), dùng cho việc tra cứu.\\
Ảo giác & Hiện tượng mô hình viết ra thông tin nghe hợp lý nhưng không có trong dữ liệu hoặc sai sự thật.\\
Nhận định & Một câu khẳng định trong câu trả lời, có thể kiểm tra đúng sai.\\
Trích dẫn & Số thứ tự gắn sau nhận định, chỉ tới nguồn làm căn cứ cho nhận định đó.\\
Token & Đơn vị văn bản mà mô hình xử lý (xấp xỉ một âm tiết hoặc một phần từ). Chi phí dùng mô hình được tính theo số token.\\
Hội thoại, lượt hỏi & Một hội thoại gồm một hoặc vài lượt hỏi nối tiếp nhau của cùng một người dùng.\\
Mô hình giám khảo & Một mô hình ngôn ngữ khác được dùng để chấm câu trả lời theo đáp án chuẩn.\\
Tập giữ kín & Bộ câu hỏi không dùng trong lúc phát triển, chỉ dùng để đánh giá.\\
Nhóm tác phẩm & Các hội thoại cùng hỏi về một tác phẩm được gộp thành một nhóm khi tính thống kê.\\
Khoảng tin cậy 95\% & Khoảng giá trị mà chênh lệch thật nhiều khả năng nằm trong đó. Nếu khoảng này chứa số 0 thì chưa thể nói hai bên khác nhau.\\
\pp{} & Hiệu của hai tỉ lệ phần trăm. Ví dụ từ 40\% lên 45\% là tăng 5 \pp.\\
API & Giao diện lập trình: cách để chương trình khác (ví dụ trang web) gửi yêu cầu tới hệ thống.\\
\bottomrule
\end{longtable}
}
